\subsection{Finite reconstruction of the actual three-division extension}

\begin{result}{A.4.11}{Proved}{An explicit characteristic refinement makes every finite deck normalizer descend to the required surface quotient}
Retain the based free group $\Sigma=\Pi_{1,1}$, its
characteristic index-nine subgroup $K_3=\Pi_{1,9}$, and
the characteristic finite quotient $\mathcal P$ of
A.4.10.  Write
$H=\Sigma/K_3\cong C_3^2$.
There is an explicitly specified characteristic finite
quotient $Q_1$ of $K_3$, and a finite extension
\[
 1\longrightarrow Q_1\longrightarrow E_2
        \longrightarrow H\longrightarrow1,
                                                        \tag{DR.1}
\]
obtained from the same quotient of $\Sigma$.
The induced deck subgroup $H\hookrightarrow
\operatorname{Out}(Q_1)$ is faithful.
The finite outer normalizer of this specified subgroup
has an actual homomorphism
\[
 N_{\operatorname{Out}(Q_1)}(H)
        \longrightarrow\operatorname{Out}(\mathcal P).
                                                        \tag{DR.2}
\]
It agrees with the ordinary descent of every geometric
normalizer through the actual three-division covering.
At the automorphism level it also gives a compatible
automorphism of a specified characteristic refinement
$P_0\twoheadrightarrow\mathcal P$, with the correct
restriction on its index-nine subgroup.

All groups and maps are finite diagonal or graph-homology
constructions.  Explicit bounds and the remaining
upper extension cocycle are given below.
The assertion concerns the specified finite normalizer;
it does not supply a geometric conjugator for a different
subgroup conjugate to $H$ in the larger outer
automorphism group.
\end{result}

\begin{proof}
For a finitely generated group $G$ and an integer $s$, put
\[
 \kappa_s(G)=\bigcap_{f:G\to\mathfrak S_s}\ker f.
                                                        \tag{DR.3}
\]
For a free group of rank $n$ this is the kernel of the
literal diagonal map specified by all $(s!)^n$ choices
of generator images.  The image has order at most
$(s!)^{(s!)^n}$.
The construction is characteristic in both the discrete
and full profinite free group, since precomposition
permutes the entire finite set of maps.

Use the explicit bound from A.4.10,
\[
 m=4(r!)^{(r!)^5},\qquad
 B_m=(m!)^{(m!)^2},\qquad q=B_m3^{B_m+1}.
                                                        \tag{DR.4}
\]
Put
\[
 N_0=\kappa_m(\Sigma),\quad P_0=\Sigma/N_0,
 \quad N_1=N_0'N_0^3,\quad E_1=\Sigma/N_1,
 \quad K_1=K_3/N_1.
                                                        \tag{DR.5}
\]
The left regular representation of $\mathcal P$, enlarged
by fixed letters to degree $m$, shows that $P_0$ maps
onto the actual quotient $\mathcal P$.
Also $m\geq9$, so $N_0\subseteq K_3$.
Writing $p=|P_0|$, the free group $N_0$ has rank $p+1$;
therefore
$|E_1|=p3^{p+1}\leq q$ and $[E_1:K_1]=9$.
The map $E_1\to P_0$ is the actual first homology cover
of the finite $P_0$-cover of the two-loop rose.

\paragraph{Faithfulness of the deck action in characteristic three.}
For every nonidentity $g\in P_0$, the rational trace of
its action on the first homology of that graph is one:
the traces on vertices and edges are zero, and the
trace on $H_0$ is one.
The identity trace is $p+1$.  Thus every nonidentity
deck element acts nontrivially on integral homology.
Its integral homology matrix has finite order, since
a power of any lift belongs to $N_0$ and acts by an
inner automorphism of $N_0$.

A finite-order integral matrix congruent to $I$ modulo
three is $I$.  Here is the needed elementary proof.
For a nonidentity such matrix write over $\mathbb Z_3$
$A=I+3^vB$, where $v\geq1$ and $B\not\equiv0\pmod3$.
Raising to an exponent prime to three leaves the
valuation $v$ of $A-I$ unchanged.
The binomial expansion gives
\[
 A^3=I+3^{v+1}B\pmod {3^{v+2}},
                                                        \tag{DR.6}
\]
so each further cubing increases this finite valuation
by exactly one.  No positive power can be the identity.
The contradiction proves the assertion.
It follows that the $P_0$-action on
$N_0/(N_0'N_0^3)$ is faithful, even when $3$ divides
$|P_0|$.
In particular
\[
 C_{E_1}(K_1)\longrightarrow P_0
           \quad\text{has image }1,                  \tag{DR.7}
\]
because a centralizer element acts trivially on the
displayed normal homology subgroup of $K_1$.

\paragraph{A characteristic nine-point input.}
Define
\[
 M=\kappa_q(K_3),\qquad Q_1=K_3/M,
 \qquad E_2=\Sigma/M.                                 \tag{DR.8}
\]
The left regular representation of $K_1$, enlarged to
degree $q$, is among the maps in (DR.3).
Consequently $M\subseteq N_1$, giving the actual
surjections $E_2\to E_1\to P_0\to\mathcal P$.
The subgroup $M$ is characteristic in the free rank-ten
group $K_3$, and $K_3$ is characteristic in $\Sigma$;
hence it is characteristic in $\Sigma$ as well.
This proves (DR.1), and gives the fully explicit bounds
\[
 |Q_1|\leq(q!)^{(q!)^{10}},\qquad
 |E_2|=9|Q_1|.                                        \tag{DR.9}
\]

Let $C=C_{E_2}(Q_1)$ and let $\mathcal E$ be the image
of the actual conjugation map
$E_2\to\operatorname{Aut}(Q_1)$.
Projection to $E_1$ and (DR.7) give
\[
 C\subseteq N_0/M\subseteq Q_1,
 \qquad C=Z(Q_1),\qquad \mathcal E=E_2/C.
                                                        \tag{DR.10}
\]
Here $C$ is abelian but is not being declared central
in $E_2$: its actual deck action is retained.
The image of $Q_1$ in $\mathcal E$ is exactly
$\operatorname{Inn}(Q_1)$, and
$\mathcal E/\operatorname{Inn}(Q_1)=H$.
Thus $H$ embeds in $\operatorname{Out}(Q_1)$, and
$\mathcal E$ is exactly its full inverse image under
$\operatorname{Aut}(Q_1)\to\operatorname{Out}(Q_1)$.
Indeed each automorphism in that inverse image differs
from one of the nine displayed deck lifts by an inner
automorphism of $Q_1$, already contained in $\mathcal E$.

Every homomorphism $\Sigma\to\mathfrak S_m$ kills
$N_0$ and therefore factors through $\mathcal E$,
using (DR.10).  Conversely each homomorphism from
$\mathcal E$ gives such a homomorphism from $\Sigma$.
The complete sets of maps are the same, and hence
\[
 \kappa_m(\mathcal E)=(N_0/M)/C,
 \qquad \mathcal E/\kappa_m(\mathcal E)=P_0.
                                                        \tag{DR.11}
\]
This identifies the lower quotient intrinsically in
the finite inner-action normalizer; it is the actual
quotient induced by $\Sigma\to P_0$.

An automorphism $\sigma$ of $Q_1$ normalizing the
specified outer subgroup $H$ conjugates its full
preimage $\mathcal E$ to itself.
It preserves the characteristic subgroup in (DR.11),
so it induces an actual automorphism of $P_0$.
Since $C=Z(Q_1)$ is characteristic, its restriction
also preserves the inverse image $N_0/M$ inside
$Q_1$ and induces exactly the same automorphism on
$K_3/N_0$.  The induced automorphism on the quotient
$H$ is the one by which $\sigma$ normalizes the deck
subgroup.  This proves compatibility of all maps in
the actual index-nine extension.

We check descent to the particular $\mathcal P$,
rather than silently treating every characteristic
free-group quotient as characteristic in $P_0$.
The subgroup $K_2/N_0$ in $P_0$ is its kernel under
the maximal elementary abelian two-quotient:
$P_0$ retains every homomorphism $\Sigma\to C_2$.
It is therefore characteristic in $P_0$.
The kernel of A.4.10 was
$N=\kappa_r(K_2)$.
Because $N_0\subseteq N$, every homomorphism
$K_2\to\mathfrak S_r$ factors through $K_2/N_0$,
and conversely.  Consequently
\[
 N/N_0=\kappa_r(K_2/N_0)
                                                        \tag{DR.12}
\]
is characteristic in that characteristic subgroup,
and hence in $P_0$.
The induced automorphism of $P_0$ therefore descends
to the exact quotient $P_0/(N/N_0)=\mathcal P$.

Replacing $\sigma$ by an inner automorphism of $Q_1$
changes the descended automorphism only by the
corresponding inner automorphism of $\mathcal P$.
This proves the well-defined homomorphism (DR.2).
For an actual geometric normalizer, all of these
conjugation maps come by restricting its action on
the original free extension $K_3\leq\Sigma$.
Their quotients are therefore exactly its ordinary
three-division descent, as asserted.
\end{proof}

\paragraph{The retained upper extension-class equation.}
The theorem proves descent to $P_0$ and $\mathcal P$;
it also records precisely the extra information
carried by the upper extension (DR.1).
Choose actual lifts $v_h\in E_2$ for the nine elements
$h\in H$, with $v_1=1$, and put
\[
 a_h=\operatorname{Ad}(v_h)|_{Q_1},\qquad
 c(h,k)=v_hv_kv_{hk}^{-1}\in Q_1.                       \tag{DR.13}
\]
A finite normalizer $\sigma$ induces some
$\nu\in\operatorname{Aut}(H)$.
For every $h$ choose an actual $q_h\in Q_1$ satisfying
\[
 \sigma a_h\sigma^{-1}
       =\operatorname{Inn}(q_h)a_{\nu h}.
                                                        \tag{DR.14}
\]
Such choices exist by the definition of the outer
normalizer, and any two choices differ by an element
of the retained group $C=Z(Q_1)$.
The discrepancy between the two products
\[
 \sigma(c(h,k))q_{hk},\qquad
 q_h a_{\nu h}(q_k)c(\nu h,\nu k)                     \tag{DR.15}
\]
lies in $C$: the two products induce the same inner
automorphism of $Q_1$ by (DR.13)--(DR.14).
Associativity on triples $(h,k,l)$ is exactly the
two-cocycle identity, with the actual $H$-action on
$C$ transported by $\nu$.
Changing the $q_h$ by a $C$-valued one-cochain changes
this discrepancy by its coboundary.
Thus (DR.15) gives the explicit finite class in
$H^2(H,C)$ governing the upper extension, together
with all its marked cochain equations.
Solving those equations is equivalent to making
the assignments $x\mapsto\sigma(x)$ on $Q_1$ and
$v_h\mapsto q_hv_{\nu h}$ into an automorphism of
$E_2$: the equations check every defining product,
and the maps on the normal subgroup and quotient
are both bijective.
No assertion that this class vanishes is used in
(DR.2).  Its coefficients die in $P_0$ by (DR.10),
which is precisely why the lower descent is already
defined.

\paragraph{The remaining finite geometric conjugacy problem.}
Conjugating $H$ by an arbitrary peripheral
automorphism of $Q_1$ transports $\mathcal E$ to
the full preimage of the resulting subgroup $H'$.
It transports (DR.11) to an isomorphic marked
quotient for $H'$.
The theorem supplies the descent after a comparison
normalizes the specified $H$; it does not identify
that transported subgroup with $H$ by a mapping
class.  The unproved step for the full A.4 census
is an actual finite geometric conjugacy comparison
between $H'$ and $H$ at this refinement, together
with the four framed insertion maps.
This locates the remaining issue in a specified
finite group with specified extension data, instead
of suppressing it under an abstract recovery of
the deck extension.
